\documentclass{article}
\usepackage{hyperref}
\usepackage{amsmath}
\usepackage{amssymb}
\usepackage[a4paper, margin=0.8in]{geometry}
\setlength{\parindent}{0pt}
\setlength{\parskip}{1em}

\begin{document}

\section{Interesting}
\begin{itemize}
\item Why do neural networks work so well?
  \item Mathematical intuition behind why different techniques work so well?
  \item What are the techniques for visualizing a model's high dimensional parameter space in 3d?
  \item Real reason for double descent?
\item What if hyperparamters weren't so fickle?
  \item Are there efficient algorithms for searching for optimal hyperparameters?
\item How is model training and inference implemented and optimized on a cluster GPUs and/or TPUs?
\item How is knowledge distilled from massive models into smaller models without a loss in performance?
\item What if we drop the distinction between training and inference and shift towards a continual learning approach?
\item What is Xavier initialization or He initialization or Layer Normalization or another init schemes supposed to be doing mathmatically?
\item What is GELU? Understand the CDF to better intuit where the percentages on the z-score table are coming from.
\item Why Warmup the Learning Rate? Underlying Mechanisms and Improvements
\end{itemize}

\section{Math concepts}
$Pr(x)$, a \href{https://medium.com/@kavya8a/understanding-probability-density-functions-a-beginners-guide-06c9821ed5c0}{probability density function}, assigns a non-negative probability density to every value in the input domain. The normalization condition is $\int Pr(x)\,dx = 1$. We assume random variables are continuous.

A joint probability distribution, $Pr(x, y)$ (probability of x and y) assigns a probability to every combination of values and $\int \int Pr(x, y) \,dxdy = 1$.

A marginal distribution is the probability of a variable while completely ignoring other variables.
\[ Pr(\mathbf{X} = x) = \sum_{y} Pr(X = x, Y = y) \]

A conditional probability $Pr(x|y)$ is the probability of x given y.
\[ Pr(x|y) = \frac{Pr(x, y)}{Pr(y)} \]

\href{https://www.datacamp.com/tutorial/kl-divergence}{KL Divergence} is a measure of how much the test distribution diverges from the correct distribution. Remember that entropy is just a weighted sum of surprisal. We subtract cross entropy from entropy. Ideally, the two distributions $p(x)$ (actual) and $q(x)$ (test) are equal, so the entropy and the cross entropy are equal, and the divergence is zero.
\[ D_{KL}[P || Q] = \int_{-\infty}^{\infty} p(x)\log{\frac{P(x)}{Q(x)}} \]

A Jacobian matrix is a matrix storing the derivative of one vector with respect to another. For example:
\[ f(x, y) = \begin{bmatrix} x + y \\ xy \end{bmatrix} \]
\[ J = \frac{\partial f}{\partial (x, y)} = \begin{bmatrix} \frac{\partial f_1}{\partial x} & \frac{\partial f_1}{\partial y} \\ \frac{\partial f_2}{\partial x} & \frac{\partial f_2}{\partial y} \end{bmatrix} = \begin{bmatrix} 1 & 1 \\ y & x \end{bmatrix} \]

The sigmoid function maps a real value $z$ to a range of $[0, 1]$.
\[ \text{sig}[z] = \frac{1}{1 + \text{exp}[-z]} \]

The softmax function maps a vector of real values (logits) $\mathbf{z}$ to a vector of probabilities in a range of 0 to 1, where all probabilities sum to 1.
\[ \text{softmax}[\mathbf{z_k}] = \frac{\exp[z_k]}{\sum_{k' = 0}^{K} \exp[z_k']} \]

Suppose a transformation $t[x]$ is applied on a function $f[x]$. The function is \textbf{invariant} if $f[t[x]] = f[x]$ and is \textbf{equivariant} if $f[t[x]] = t[f[x]]$.

The mean of a distribution is it's balancing point. The summed distance of all the values on the left of the mean and all the values on the right of the mean will cancel out. Can think of the distribution as a weight scale, where the mean is the pivot. Expectation is the mean. It's defined as the weighted sum of each value in the distribution by the probability (can think of it from a frequentist point of view) of that value.

\section{Neural network basics}
A model is a mathematical function defined by parameters that maps an input to an output. There are discriminative models and generative models. Discriminative models find outputs from inputs, while generative models find inputs from outputs. Discriminative models exhibit \href{https://en.wikipedia.org/wiki/Transduction_(machine_learning)}{transductive} reasoning, mapping specific inputs to specific outputs. Generative models exhibit inductive reasoning, mapping specific inputs to generalized rules. Most models are discriminative since it's easier to exploit the prior knowledge that's baked into the input data.

Using a non linear activation function allows a model to represent non linear functions. A saturated activation function is one where its derivative approaches zero when the magnitude of its output becomes very large (very positive or very negative). Using a saturated activation function would hamper training as gradients approach zero in deeper and deeper layers (vanishing gradient problem). A common activation function that is non saturated is ReLU (Rectified Linear Unit) and Leaky ReLU. ReLU makes the activations in a layer sparse, while leaky ReLU preserves a small gradient for negative activations, which prevents those activations from dying.
\[ \text{ReLU}(x) = \begin{cases} x & \text{if } x \geq 0 \\ 0 & \text{if } x < 0 \end{cases} \]
\[ \text{Leaky ReLU}(x) = \begin{cases} x & \text{if } x \geq 0 \\ \alpha x & \text{if } x < 0 \end{cases} \]

Network capacity is the model's ability approximate different functions. The network capacity is defined as the number of parameters, the number of layers (depth), the number of hidden units in each layer, etc. The theoretical capacity is called the representational capacity, while the actual number of functions that the model can approximate is called the effective capacity. With enough capacity (hidden units), the network can approximate any function at a given resolution (Universal approximation theorem).

Parameters are layer weights and biases, while hyperparameters are high level settings on how the neural network operates (ex: learning rate, epoch, batch size, number of layers, hidden units per layer, etc). A common way to optimize hyperparameters is to split the training dataset into three subsets. A training set, a validation set, which is used to choose the hyperparameters that lead to the best model performance, and a test set which is used for a final model performance evaluation. Another common approach is called k-fold cross validation. Training samples are split into $N$ random subsets, and for $N$ iterations, training is done with $N - 1$ subsets while hyperparameter validation is done with the remaining subset. The final model performance is evaluated using an average of the $N$ model predictions on a different test dataset.

There are three main sources of error when training a model: The extent to which model performance is maintained on the training set is called generalization.
\begin{itemize}
\item \textbf{Noise:} The inherent randomness (mislabelling, multiple inputs mapping to the same output, etc) in the training dataset is insurmountable.
\item \textbf{Bias:} The model might not have enough capacity to fit the data accurately. Increasing the model capacity decreases bias but increases variance, and doesn't necessarily decrease test error. This is known as the bias-variance tradeoff. Increasing model capacity unfortunately makes the model fit noise in the training data better, which is referred to as overfitting.
\item \textbf{Variance:} The fitted function will vary slightly based off of the training samples (and possibly SGD). Increasing the number of training samples decreases variance.
\end{itemize}

As model capacity increases, test error should decrease up until the point where variance becomes too large and the model starts to overfit, making the test error increase. However, for many datasets, the test error starts decreasing again as the model capacity continues to increase. This is referred to as \href{https://www.lesswrong.com/posts/FRv7ryoqtvSuqBxuT/understanding-deep-double-descent}{double descent}. At the moment we don't fully know why, but there's at least one plausible explanation. As the model capacity increases and the training error approaches zero, the model's fitting function is able to pass through every single training sample. The segments between data points are initially very erratic, leading to a large test error. However, as the model capacity continues to grow, the segments between data points become smoother and smoother, which makes the model fit to new test points better, decreasing test error. This is especially plausible when you consider that in a high dimensional input space, the training samples are extremely sparse (curse of high dimensionality). There's a large number of possible fitting functions passing through all training samples that the model can adopt. The model's tendency to adopt certain families of fitting functions is called inductive bias. Any factor that biases a solution to a set of equivalent solutions is called a regularizer. The hypothesis is that SGD is an implicit regularizer of smooth functions.

If the expectation of the magnitude of the weights are smaller than 1, the gradient vanishes, and if the expectation of the magnitude of the weights are larger than 1, the gradient explodes. Both problems can occur in the forward or backward pass. A solution to this is He initialization (also known as Kaiming initialization). He initialization keeps the variance of layer activations relatively stable. Initializing the weights to random values taken from a normal distribution centered around 0, with a variance (1), where $D_h$ is the dimension of the current layer, and $D_{h'}$ is the dimension of the next layer, keeps the gradients stable. This particular initialization is used with the ReLU activation function. If the current and next layer have the same dimension, then an average of the variance for the forward pass and backward pass doesn't have to be used (2).
\begin{align*}
\sigma^2 &= \frac{4}{D_h + D_{h'}}\tag{1}\\
\sigma^2 &= \frac{2}{D_h}\tag{2}
\end{align*}

\section{Loss functions}
A model outputs a set of parameters ($\phi$) that describe a probability distribution. Minimizing the loss function $L[\phi]$ means maximizing $Pr(y|\phi)$. During inference, the final output value is the maximum of the output probability distribution.
\begin{align*}
\boldsymbol{\hat{\phi}} &= \underset{\phi}{\text{argmax}}\left[\prod_{i = 1}^{I} Pr(y_i|x_i)\right]\\
                        &= \underset{\phi}{\text{argmax}}\left[\prod_{i = 1}^{I} Pr(y_i|\theta_i)\right]\\
                        &= \underset{\phi}{\text{argmax}}\left[\prod_{i = 1}^{I} Pr(y_i|f[\mathbf{x_i}, \boldsymbol{\phi}])\right]
\end{align*}

It's more practical to take the sum of the log of all the probabilities, since they may be very small. Also, by convention the loss function should minimize loss, which can be achieved by multiplying the sum by -1 and using argmin. This is the negative log-likelihood!
\[ \boldsymbol{\hat{\phi}} = \underset{\phi}{\text{argmin}}\left[-\sum_{i = 1}^{I} \log[Pr(y_i|f[\mathbf{x_i}, \boldsymbol{\phi}])]\right] \]

The following are a few common loss functions.
\begin{itemize}
\item \textbf{Mean squared error} is used in regression tasks, when the output follows a Gaussian distribution. When performing inference, the model's raw output can be used since it is the mean of the output gaussian distribution, which so happens to be the distribution's maximum. Normalization is done to make the loss independent of the number of training samples, and to make comparing loss across datasets easier. $I$ is the number of training samples.
\[ L[\phi] = \frac{1}{I} \sum_{i = 1}^{I} (\mathbf{y_i} - f[\mathbf{x_i}, \boldsymbol{\phi}])^2 \]

\item \textbf{Cross entropy loss} is used in classification tasks, when the output follows a categorical distribution. A categorical distribution is a distribution where the probability of each class is $\lambda_k \in [0, 1]$ and $\sum_{K} \lambda_k = 1$. To ensure that the model output follows a categorical distribution, softmax is applied. $y$ is the true class, $\hat{y_i}$ is the predicted class.
\[ L[\phi] = -\sum_{i = 1}^{I}\sum_{k=1}^{K} y_i \log(\hat{y_i}) \]

\item \textbf{Dice loss} is often used with U-Net models. It's how many values are in both samples $P$ and $Q$ divded by the total number of samples in $P$ and $Q$.
\[ \frac{|P \cup Q|}{|P| + |Q|} \]

\item When the model's output is multivariate, each output and each error are treated as independent. The model's loss function becomes a sum of all the loss functions for each output. See the chart on p. 70 of textbook for more.
\end{itemize}

\section{Backpropagation}
Backpropagation is an algorithm that iteratively applies the chain rule to any arbitrary computational graph. A neural network can be thought of as a series of function compositions, where each layer is a function that composes the previous layer. A neural network can also be thought of as a \href{https://www.youtube.com/watch?v=i94OvYb6noo}{directed acyclic graph}, where each layer is a node. Backpropagation is just going through each node, right to left, and computing the local gradient, multiplying it by the global gradient (which is initialized to 1), and passing an updated global gradient to child nodes (previous layers). Gradients are computed for every training sample in a batch, then summed together to get the gradient used in an optimization algorithm. The loss function gradient is used as the initial gradient.

A forward pass is done to compute and cache values of hidden units in each layer, which will be used to compute gradients. Each sample in the batch uses the same weights and the same biases.
Layers are different for every sample in the batch, and weight and bias gradients are summed across the batch (don't have a batch dimension).

\begin{align*}
f_0       &= \boldsymbol{\beta_0} + \boldsymbol{\Omega_0}\mathbf{x}\\
h_k       &= a[f_{k - 1}] && k \in \{1, 2, 3, \dots, K\}\\
f_{k - 1} &= \boldsymbol{\beta_k} + \boldsymbol{\Omega_k} \mathbf{h_k} && k \in \{1, 2, 3, \dots, K\}
\end{align*}

A backward pass is done to compute the derivative of the loss with respect to every weight and bias in the network. The rules for computing local gradients and updating the global gradient are as follows. Note that an outer product is used when computing the derivative of loss with respect to weights. The derivative of the loss with respect to the hidden units is a column vector. Gradients are layer specific, and during the optimization algorithm, the gradient used to step is an average of all the gradients for each training sample in the batch.
\[ \frac{\partial l_i}{\partial \boldsymbol{\beta_k}} = \frac{\partial l_i}{\partial f_k} \]
\[ \frac{\partial l_i}{\partial \boldsymbol{\Omega_k}} = \frac{\partial l_i}{\partial f_k} h_k^T \]
\[ \frac{\partial l_i}{\partial f_{k - 1}} = I[f_{k - 1} > 0] \odot \boldsymbol{\Omega_k}^T \frac{\partial l_i}{\partial f_k} \]

$k \in \{K - 1, K - 2, K - 3, \dots, 1\}$, and $\frac{\partial l_i}{\partial f_k}$ is the global gradient. The derivative of ReLU with respect to the hidden units is a Jacobian matrix ($J$). When $i \neq j, J_{ij} = 0$, since $\text{ReLU}^{\prime}(i)$ only depends on $i$. The diagonal elements in $J$ would be either 1 or 0 (considering the definition of ReLU). Instead of computing the entire Jacobian, its diagonal elements are directly computed, which is just: \texttt{hidden\_unit > 0}.

Backpropagation is different for \href{https://deeplearning.cs.cmu.edu/F21/document/recitation/Recitation5/CNN_Backprop_Recitation_5_F21.pdf}{convolutional layers}. We need to compute the derivative of the loss with respect to the kernel F (1) and the derivative of the loss with respect to the previous layer X (2).
(1) X is padded by the current layer's padding and the gradient is dilated by the current layer's stride.
(2) To rotate by 180 degrees, F is flipped vertically then horizontally. The gradient is dilated by the current layer's stride and zero-padded by the current layer's kernel size minus one. After the convolution, the gradient should have its padding stripped.

\begin{align*}
\frac{\partial l_i}{\partial F} &= \text{convolution}\left[X, \frac{\partial l_i}{\partial f_k}\right] \tag{1} \\\\
\frac{\partial l_i}{\partial X} &= \text{convolution}\left[\frac{\partial l_i}{\partial f_k}, \text{rot180}[F]\right] \tag{2}\\
\frac{\partial l_i}{\partial B_k} &= \sum_{y=1}^{H} \sum_{x=1}^{W} \frac{\partial l_i}{\partial f_{k,y,x}} \tag{3}
\end{align*}

\section{Optimization algorithms}
The goal of an optimization algorithm is to update the model's parameters as to minimize the loss. This process is known as learning, training or fitting.

\begin{itemize}
\item \textbf{Gradient descent:} Gradients are the partial derivatives of the loss function with respect to the model parameters. $\alpha$ is the learning rate, which is a hyperparameter that controls how quickly the parameters change. A \href{https://machinelearningmastery.com/a-gentle-introduction-to-learning-rate-schedulers/}{learning rate} scheduler like StepLR is often used ($\gamma$ is a decay multiplier). The problem with gradient descent is that there's no way of knowing whether the model will converge on a global minimum or a local minimum, or get stuck in a saddle point (flat area of the function) especially when the model function is non-convex.
\[ \alpha_t = \alpha_0 * \gamma^{\frac{\text{epochs}}{\text{step\_size}}} \]
\[ \boldsymbol{\hat{\phi}} = \boldsymbol{\phi} - \alpha\frac{\partial L}{\partial \boldsymbol{\phi}} \]

\item \textbf{Stochastic gradient descent:} SGD adds randomness to the stepping process by sampling a random subset of training samples (batch) every epoch. Because of the batching, SGD is more efficient. Since the loss function for each batch is different, gradients vary, which makes the model take a noisier path to convergence. This allows the parameters to potentially step out of a local minima or a saddle point.
\[ \hat{\boldsymbol{\phi}} = \boldsymbol{\phi} - \alpha \sum_{i \in B_t} \frac{\partial l_i[\boldsymbol{\phi_i}]}{\partial \boldsymbol{\phi}} \]
Where $B_t$ is the current batch of training samples and $l_i$ is the loss given the chosen training samples. SGD generalizes better than gradient descent, and smaller batch sizes result in better performance. This is because SGD implicitly penalizes variance, making the model learn a function that's a tighter, smoother fit to the data.

\item \textbf{Nesterov accelerated momentum:} By adding a weighted combination of the previous gradients, and the gradients at the target position, the model can converge on a minimum much faster. Notice that the momentum terms are recursive, and older momenta get weighed less and less. Past gradients are weighed more than the current gradients. By computing the gradient at the target position we also avoid the overshooting that would happen if we used the gradients at the starting position.
\[ m_{t + 1} = B * m_t + (1 - B)\sum_{i \in B_t} \frac{\partial l_i[\boldsymbol{\phi_t} - \alpha B * m_t]}{\partial \boldsymbol{\phi}} \]
\[ \boldsymbol{\hat{\phi}} = \boldsymbol{\phi_t} - \alpha * m_{t + 1} \]

\item \textbf{Adam:} The Adam optimizer takes ideas from several other algorithms and combines them into one. First, Include a weighted history of both the momentum and the squared momentum. Second, amplify the momentum and squared momentum when t is small, doing nothing when t is large. This is because the momentum (which is initialized to 0) will be too small when t is small. Lastly, normalize the gradients so that the parameters step the same distance in each direction. Momentum is preserved across batches and epochs.
\begin{align*}
m_{t + 1} &= \beta m_t + (1 - \beta)\sum_{i \in B_t} \frac{\partial l_i(\boldsymbol{\phi}_t)}{\partial \boldsymbol{\phi}} \\
v_{t + 1} &= \gamma v_t + (1 - \gamma)\left(\sum_{i \in B_t} \frac{\partial l_i(\boldsymbol{\phi}_t)}{\partial \boldsymbol{\phi}}\right)^2 \tag{1} \\
\tilde{m}_{t + 1} &= \frac{m_{t + 1}}{1 - \beta^{t + 1}} \\
\tilde{v}_{t + 1} &= \frac{v_{t + 1}}{1 - \gamma^{t + 1}} \tag{2} \\
\boldsymbol{\hat{\phi}} &= \boldsymbol{\phi}_t - \alpha \frac{\tilde{m}_{t + 1}}{\sqrt{\tilde{v}_{t + 1}} + \epsilon} \tag{3}
\end{align*}
\end{itemize}

\section{Regularization}
See p.155, figure 9.14

Regularization refers to any method that decreases the generalization gap between test and training performance. Explicit regularization involves adding an explicit term to the loss function in order to favor certain parameters. This involves moving away from a maximum likelihood criterion to a maximum a posteriori criterion, where the negative log loss also includes a probability $Pr(\phi)$, a belief about the data.
\[ \boldsymbol{\hat{\phi}} = \underset{\phi}{\text{argmax}}\left[Pr(\phi)\prod_{i = 1}^{I} Pr(y_i|x_i, \phi)\right] \]

The most common regularization term is the \textbf{L2 norm}, which biases smaller parameters, making the model output change slower. Parameters should be as small as possible (simplest solution) while still minimizing the loss function. L2 norm is applied to weights, and having smaller weights will reduce variance and make the model fit a smoother function, at the cost of reducing data fit. With small values of $\lambda$, test performance improves.
\[ \boldsymbol{\hat{\phi}} = \underset{\phi}{\text{argmin}}\left[\sum_{i = 1}^{I} l_i[x_i, y_i] + \lambda \sum_{j}\phi_j^2\right] \]

The following are ways to improve model performance. Note that these aren't silver bullets. In theory they should improve test performance, but neural networks are very finicky; the optimal hyperparameters and learning strategies to use should be determined empirically.
\begin{itemize}
  \item \textbf{Early stopping}: Stop training the model before it fully converges. In practice this means training the model for a large number of epochs to see which number of epochs maximizes performance.
  \item \textbf{Ensembling}: Run inference on multiple models and average their outputs together to get a more confident answer. This can be done by training identical models with random parameter initializations, identical models trained on different randomly sampled training data (called bootstrap aggregating), or training different model architectures on identical data.
  \item \textbf{Dropout}: During training, clamp a random subset of hidden units (typically half) to zero. This helps the model reduce its reliance on any given hidden unit and encourages the weights to be smaller so that the output fluctuates less wildly when hidden units are turned off. During inference, all hidden units are active, so weights are multiplied by $1 - \eta$, where $\eta$ is the dropout probability (weight scaling inference rule). One can also do Monte Carlo dropout instead, where you run inference on the model multiple times with random subsets of the hidden units clipped, and then combine results.
  \item \textbf{Apply noise} to:
  \begin{itemize}
    \item Inputs: Generate \href{https://arxiv.org/html/2502.18691v1}{different variations of the input data} in order to teach the model to be indifferent to irrelevant data transformations. These are referred to as data augmentation and the augmentation techniques are applied per each minibatch on the fly. Makes the model fit to more variable samples, improving test performance.
    \item Weights: Makes the training converge to a local minima inside of a flat middle region, where changing the weights doesn't affect the output much. The model becomes less particular about specific weight values.
    \item Labels: Makes the model stop chasing absolute certainty. Label smoothing refers to changing the loss function to minimize the cross-entropy between the predicted and actual distribution (Where the true label has a probability $1-\rho$ and all the other classes have equal probability. $\rho$ is the proportion of the training labels that are assumed incorrect.).
  \end{itemize}
  \item \textbf{Transfer learning}: The initial layers of a model act as feature extractors and the final layers act as classification heads. Transfer learning leverages this by freezing the parameters of the pre-trained feature extractor and only training the classification head on the target dataset. Multi-task learning is when the model is trained to solve multiple problems concurrently, hoping that the performance on each will improve.
  \item \textbf{Self-supervised learning}: Generate more labelled data using a secondary generative model.
  \begin{itemize}
    \item In generative self-supervised learning, parts of each data is masked and the secondary task is to predict the missing part.
    \item In contrastive self-supervised learning, pairs of examples with commonalities are compared to unrelated pairs.
  \end{itemize}
\end{itemize}

\section{Convolutional networks}
A convolutional network is composed of a series of convolutional layers. Convolutional layers are a special subset of fully connected layers. A convolution is a weighted sum over a set of inputs using fixed parameters called a convolutional kernel or a filter. A valid convolution is a convolution that avoids inputs that are out of range, and a zero padded convolution is a convolution that uses zeros in place of inputs that are out of range. Stride is the step size used when sliding over the input and dilating is the number of zeros interspersed between the inputs plus one.

A convolutional layer computes its outputs by convolving the input, adding a bias and passing the result through an activation function. Each kernel used to convolve corresponds to an output channel. Channels add depth to the convolutional layer. A convolutional layer is equivariant to translation because of the use of the same kernel over the entire image, and partially invariant to translation because of max pooling.

For a 2D convolution on RGB images, a $C_i \times K \times K$ kernel is slid over the image, resulting in a $C_{h} \times C_{w}$ matrix. $C_i$ is the number of input channels, $C_{w}$ is the output width, $C_{h}$ is the output height. Hidden layers are 3D, unlike layers in feedforward networks that are always 1D. Each channel's hidden unit is computed as the following, where $\omega_{ij}$ is the kernel weights, $m$ is $\lceil K/2 \rceil$ and $n \in [1, C_o]$.
  \[ h_{nyx} = a\left[\boldsymbol{\beta} + \sum_{c=1}^{C_i} \sum_{i=1}^{K} \sum_{j = 1}^{K} \omega_{cij} x_{c,i-m,j-m}\right] \]

Weights are defined as $\boldsymbol{\Omega} \in \mathbb{R}^{C_o \times C_i \times K \times K}$ where $C_i$ is the number of channels in the previous layer, $C_o$ is the number of channels in the current layer and $K$ is the kernel size. Biases can be defined as $\boldsymbol{\beta} \in \mathbb{R}^{C_o}$. The output layer size is defined as the following where $N$ is the current layer size, $P$ is padding and $K$ is the kernel size and $S$ is the stride.
\[ \left\lfloor \frac{N + 2P - K}{S} + 1 \right\rfloor \]

There are different types of convolutions:
\begin{itemize}
  \item \textbf{Strided convolution}: Refers to a convolution with a stride $> 1$. This downsamples the original input size. For example, with a stride of 2, there would be half as many outputs as there are inputs. There are other ways to downsample an image. Using max pooling/mean pooling reduces a 2x2 block of pixels down to the max value/mean value of the block.
  \item \textbf{Transposed convolution}: Refers to a convolution that uses the transpose of the weight matrix used in a strided convolution. This upsamples the original input size. For example, there would be twice as many outputs as there are inputs, and each input would correspond to two outputs. There are other ways to upsample an image. Using max unpooling or bilinear interpolation or duplicating the current value across a 2x2 block.
  \item A \textbf{1x1 convolution} reduces the number of feature maps without reducing the number of spatial dimensions. A $C_o \times C_i \times 1 \times 1$ kernel is ran on the input layer, resulting in a $1 \times C_h \times C_w$ result for every output channel $C_o$. Since this mixes channels this also adds non linearity.
\end{itemize}

\section{Normalization}
BatchNorm is a technique that helps optimize the training process. It consists of a layer with learnable parameters $\gamma$ and $\delta$, and running averages of the parameters. The layer first standardizes the input before applying the linear combination.
The mean and the variance are summed across the batch dimension (linear dimension $m \times D \rightarrow D \times 1$, convolutional dimension $m \times C_i \times H \times W \rightarrow C_i$). During inference the batch size is 1, which makes the resulting variance = 0, breaking standardization. Instead of standardizing then applying a linear combination, only a linear combination with the rolling averages of the parameters is done.
The dimensions of $\gamma$ and $\delta$ in a linear layer is $D \times 1$ and in a convolutional layer is $C_i \times 1$. $m$ is the batch size, $\epsilon$ prevents division by 0.
\[ \mu = \frac{1}{m} \sum_{i \in B} x_i \]
\[ \sigma = \sqrt{\frac{1}{m} \sum_{i \in B} (x_i - \mu)^2 } \]
\[ x = \frac{x - \mu}{\sigma + \epsilon} \]
\[ y = \gamma x + \delta \]

By standardizing, the input is constrained to a smaller range of values (centered around zero, variance of one). The learnable parameters enable the input to maintain an optimal and stable numerical scale. This smoothes the model's loss landscape and enables gradient descent to take larger and more reliable steps, helping deep networks train successfully.
Because each normalized output depends on every other input, its gradient is more involved.
\[ \frac{\partial l_i}{\partial x_i} = \frac{\gamma}{m \sigma} \left( m g_i -  \sum_{j} g_j - \hat{x_i} \sum_{j} g_j \hat{x_j} \right) \]
$g_i$ is the global gradient for a specific activation $i$, $\hat{x_i}$ is the normalized activation.

LayerNorm is another normalization technique that does the same thing as BatchNorm, except the mean and the variance are summed across the layer dimension (linear dimension $m \times D \rightarrow m \times 1$, convolutional dimension $m \times C_i \times H \times W \rightarrow m$).
This changes the statistics used to normalize. Instead of normalizing a feature across multiple samples, a feature is normalized across the layer's own internal representation. LayerNorm is mostly used in RNNs and Transformers while BatchNorm is used in CNNs.

\section{Residual networks}
Training deeper and deeper networks unexpectedly results in poorer performance. One possible reason for this is that the deeper the network, the less smooth its loss landscape is, which means that the gradient doesn't remain similar between steps. In a deep network, small changes in the input result in wildly different gradients (shattered gradients problem). Additionally, small changes in earlier layers results in large and unpredictable changes in later layers.

A residual network comprises of a series of residual blocks. A residual block is a set of layers that ends in a residual connection, which is just adding the output of the previous block to the output of the current block. In a residual block, BatchNorm is applied first to normalize the inputs. ReLU comes before the linear combination, in order to preserve the gradient of the residual block (except in the first residual block in order to preserve negative values in the input). (p. 191, figure 11.6). Each block comprised of the following layers in the original paper: BatchNorm, ReLU, Convolution, BatchNorm, ReLU, Convolution.
Often, the block input size will not match the block output size. In such cases, a projection (1x1 convolution with a stride chosen to downsample/upsample the input) is used. The projection is also normalized with a corresponding BatchNorm layer.
\[ h_k = h_{k - 1} + f[x, \phi] \]
The residual network can be thought of an ensemble of smaller subnetworks summed together. It can also be thought of as a computational graph where each layer (input included) takes several different paths of varying length to the output. (p. 190)

Architectures using residual connections includes ResNet, DenseNet, \href{https://openaccess.thecvf.com/content/WACV2021/papers/Zhang_ResNet_or_DenseNet_Introducing_Dense_Shortcuts_to_ResNet_WACV_2021_paper.pdf}{Dense shortcuts for ResNet}, U-Net.

\section{Reccurent Neural Networks}
\href{https://stanford.edu/~shervine/teaching/cs-230/cheatsheet-recurrent-neural-networks/}{Reccurent Neural Networks} are models designed for processing sequences of data. Depending on how it's designed, an RNN can predict a sequence from a sequence, a sequence from a value or a value from a sequence. Here, sequence to sequence prediction will be used to \href{https://www.mdpi.com/information/information-15-00517/article_deploy/html/images/information-15-00517-g001-550.jpg}{explain how RNNs work}.
An RNN is composed of hidden blocks $h_t$ which output a symbol $o_t$ in the output sequence given an input sequence symbol $x_t$ and a running sum of the previous hidden blocks $h_{t - 1}$. $h_t$ is a compressed representation of all the symbols previously seen in the input sequence at position $t$. Each hidden block is defined as $(1)$ and each output is defined as $(2)$, where $a_h$ and $a_o$ are the activation functions for the hidden blocks and the output, $b_h$ and $b_o$ are the biases for the hidden blocks and the output, $W_x$ is the input weight, $W_h$ is the hidden-state-to-hidden-state weight and $W_o$ is the output weight.
\[ h_t = a_h(W_x x_t + b_h + W_h h_{t - 1}) \tag{1} \]
\[ o_t = a_o(W_o h_t + b_o) \tag{2} \]
$W_x$, $W_h$ and $W_o$ are shared by all the hidden blocks in the network. Thus, an RNN can be thought of as the same shallow network that remembers a window of past symbols, reused repeatedly on each symbol in the input sequence,

RNNs have several flaws. The most obvious one is that RNNs are inherently sequential and thus unparallelizable. This makes them very inefficient to train and perform inference. Also, the number of symbols an RNN can remember at once is bound by floating point precision, which makes that number limited. Lastly, the vanishing/exploding gradient problem is still an issue, especially since each hidden block shares the same set of weights.

\section{Transformers}

\href{https://arxiv.org/abs/1409.0473}{Neural Machine Translation By Jointly Learning to Align and Translate} introduced the concept of \textit{attention} to improve the way that RNN encoder-decoder models process sequences. By weighing certain input symbols more than others in the context vectors used by the decoder, the model can focus on the most relevant parts of the input sequence needed to predict the next output symbol. However, the sequential nature of RNNs prevents parallelization, and the compute needed to process a sequence grows proportionally to its length. \href{https://arxiv.org/abs/1706.03762}{Attention Is All You Need} introduces the Transformer, a model architecture that performs better than the RNN models used previously while addressing their inefficiencies.

The input sequence and previously outputted sequence are converted into a sequence of tokens. For example, \href{https://en.wikipedia.org/wiki/Byte-pair_encoding}{Byte-pair encoding} can be used to convert text into a sequence of tokens that are word fragments. These tokens are then projected into embeddings, which are points in high dimensional space. The model learns to position each point relative to the other points in a way that encodes semantic meaning. Techniques like cosine similarity or the dot product can be used to find points that are close by, and thus are semantically similar. A fixed positional embedding is added to the projected vectors to order them.

Embeddings for the input tokens are processed by a series of encoder layers. Each encoder layer first computes self-attention, where the query, keys and values are derived from the same input sequence. This is done to attend to the entire input sequence at once so that necessary context can be incorportaed into each embedding before they are passed into a feed forward network. The feed forward network is composed of two layers, $\text{ReLU}(\mathbf{x} W_1 + b_1) W_2 + b_2$, which process each embedding independently in parallel.

Embeddings for the previously outputted sequence are processed by a series of decoder layers. Each layer first computes masked self-attention, where a causal mask ($1$ below the diagonal, $-\infty$ on and above the diagonal) is multiplied to the attention logits before softmax is applied, preventing each position from attending to future tokens. This maintains the Transformer's autoregressive nature. Cross-attention is then performed, where $Q$ is from the current decoder layer and $K, V$ are from the very last encoder layer, allowing each output symbol to attend to all input symbols at once.

Each encoder and decoder layer maintains residual connections between sublayers. This is to preserve the original token identity and position all the way up to the final output.

Multi-head attention is used as the core attention mechanism. First, the input embeddings are projected into a smaller vector subspace $h$ times, producing $h$ different representations of $Q, K, V$. These operations are equivalent to projecting into the same vector space and slicing $h$ times. This allows each head to compute attention using a different learned notion of similarity. Next, each attention head computes how similar each query is to each key and stores that in a similarity matrix. The similarity matrix is scaled by $1 / \sqrt{d_k}$ which prevents the dot product from growing excessively large as $d_k$ increases, which would cause softmax to become overly peaked, resulting in tiny gradients. Each row in the similarity matrix is used to weigh how relevant each value is to the query and softmax is used to normalize the raw logits into a probability distribution. However, the use of softmax gives this attention variant a quadratic time complexity. For Attention, $Q \in \mathbb{R}^{n_q \times d_k}$, $K \in \mathbb{R}^{n_k \times d_k}$, $V \in \mathbb{R}^{n_k \times d_v}$, $M \in \mathbb{R}^{n_q \times n_k}$, where $d_k$ and $d_v$ are embedding dimensions of the keys and values.

\begin{align*}
\text{MultiHead}(Q, K, V) &= \text{Concat} \left( \text{head}_1, ..., \text{head}_h \right) W^O \\
\text{head}_i &= \text{Attention}(Q W_i^Q, K W_i^K, V W_{i}^V) \\
\text{Attention}(Q, K, V) &= \text{softmax} \left( \frac{Q K^{\top}}{\sqrt{d_k}} \odot M \right) V
\end{align*}

After the embeddings, encoder layers and decoder layers, there is a final linear layer that outputs a $n \times V$ matrix of probability distribution, where $n$ is the number of embeddings and $V$ is the vocabulary size. Each row in the matrix corresponds to a probability distribution of the following word. For example, if the input is "<BOS> I like", the output might correpond to "I like cats". "I" is predicted to come after "<BOS>", "cats" is predicted to come after "like" and so on. Thus, during training the output token sequence must be shifted right before being fed into the Transformer so that the model doesn't learn to simply copy the last token in the sequence.

\section{Generative Pre-trained Transformer models}

Instead of having an ensemble of task-specific models and datasets, it would be better to have a
single model that has zero-shot task transfer with as little parameter or architectural changes as possible.

\subsection{GPT-1}
\begin{itemize}
\item It's composed of several layers of transformer decoders followed by a feed forward head.
\item The transformer is referred to as a "decoder" because of its masked self-attention, not because of its structure.
\item Training is split into two phases: generative pre-training and supervised fine tuning.
\item Pre-training is self-supervised since the likelyhood of predicting the next token is optimized.
\item During post-training, any arbitrary supervised learning objective can be optimized.
\end{itemize}

\subsection{GPT-2}
\begin{itemize}
\item Maximize p(output | input, task), where the output, input and task are a sequence of symbols. Instead of seperating input, task and output, combine them into one input symbol sequence.
\item In principle, language modelling should be able to learn various tasks through next token prediction. Although the objective is extremely simple, the skills needed to do it effectively are very complex. For example, to predict the next word in this sentence: "In french we would say the phrase: 'The dog ate my homework' as 'Le chien a mangé mes" as "devoirs", the model needs to have a semantic and gramatical understanding of the english sentence and partial french sentence, it needs to recognize that this is a translation task, and it needs to understand the nuance between the different next word choices, etc.
\item Given enough capacity and enough diverse data, both in massive quantities, the model should be able to learn the abstract representations needed to model language effectively and generalize well to unseen examples. The question is *how*.
\item Text is tokenized using BPE (on bytes not unicode chars) and merging across punctuation is not allowed to enable better use of a limited vocabulary.
\item The model follows the original GPT-1 architecture with a few changes. The context window size is 1024 tokens long and residual layers weights are scaled by 1 / sqrt(N), where N = number of residual layers. LayerNorm as the first step in layer in order to normalize inputs and preserve an identity gradient and another LayerNorm after the final self-attention block.
\end{itemize}

\end{document}
